% from ega4-1-fr.tex
\oldpage[IV-1]{240}

\begin{proof}
On se ramène comme ci-dessus à prouver que $f(X)$ est constructible.
Appliquons le critère (0\textsubscript{III}, 9.2.3)~: il faut prouver que pour
toute partie fermée irréductible $T$ de $Y$, $T\cap f(X)=f(f^{-1}(T))$ est
rare dans $T$ ou contient un ouvert non vide de $T$~; en prenant un
sous-préschéma de $Y$ ayant $T$ pour espace sous-jacent (I, 5.2.1) et en
considérant son image réciproque par $f$, on voit finalement (compte tenu de
(1.5.4, (iii)) qu'on est ramené à prouver que si l'on suppose $Y$
\emph{irréductible et $f$ dominant}, $f(X)$ contient un ouvert non vide de
$Y$. On peut en outre supposer $Y$ affine~; alors $X$ est réunion finie
d'ouverts affines $V_i$, et comme $f(X)$ est dense dans $Y$, il en est de
même d'un au moins des $f(V_i)$. On peut donc aussi supposer $X$ affine~; si
$(X_j)$ est la famille (finie) des composantes irréductibles de $X$, on voit
comme ci-dessus qu'un au moins des $f(X_j)$ est dense dans $Y$~; on peut donc
supposer $X$ irréductible. Enfin, en remplaçant $f$ par $f_{\mathrm{red}}$
(1.5.4, (vi)) on peut supposer $X$ et $Y$ \emph{intègres}. Alors
$X=\operatorname{Spec}(B)$, $Y=\operatorname{Spec}(A)$, où $A$ est un
anneau intègre noethérien, $B$ une $A$-algèbre intègre de type fini contenant
$A$ (I, 1.2.7). Il suffit de montrer qu'il existe $g\in A$ tel que, pour tout
$y\in D(g)$ (c'est-à-dire tel que $g\notin\mathfrak p_y$) l'idéal
$\mathfrak p_y$ soit l'intersection de $A$ et d'un idéal premier de $B$, car
cela montrera que $D(g)\subset f(X)$. Finalement il suffit de prouver le
\end{proof}

\begin{lemma}[1.8.5.1]
\label{IV.1.8.5.1-fr}
Soient $A$ un anneau intègre, $B$ une $A$-algèbre intègre de type fini. Il
existe $g\in A$ tel que tout homomorphisme de $A$ dans un corps
algébriquement clos $\Omega$, non nul en $g$, se prolonge en un homomorphisme
de $B$ dans $\Omega$.
\end{lemma}

\begin{proof}
Or c'est là un résultat classique d'algèbre commutative (Bourbaki,
\emph{Alg. comm.}, chap. V, \S~3, n\textsuperscript{o}~1, cor. 3 du th. 1).
\end{proof}

\begin{corollary}[1.8.6]
\label{IV.1.8.6-fr}
Soient $Y$ un préschéma irréductible, $\eta$ son point générique,
$f:X\to Y$ un morphisme localement de type fini. Si $f^{-1}(\eta)$ n'est pas
vide, il existe un voisinage ouvert $U$ de $\eta$ dans $Y$ tel que
$f^{-1}(y)$ soit non vide pour tout $y\in U$. Si $\mathscr F$ est un
$\mathcal O_X$-Module quasi-cohérent de type fini, et si
$\mathscr F_\eta=\mathscr F\otimes_{\mathcal O_Y}k(\eta)$ n'est pas nul,
alors il existe un voisinage ouvert $U'$ de $\eta$ dans $Y$ tel que
$\mathscr F_y=\mathscr F\otimes_{\mathcal O_Y}k(y)$ ne soit pas nul pour
tout $y\in U'$.
\end{corollary}

\begin{proof}
Si $p_y$ est la projection canonique de la fibre
$f^{-1}(y)=X\times_Y\operatorname{Spec}(k(y))$ dans $X$, on a
$\mathscr F_y=p_y^*(\mathscr F)$, donc
$\operatorname{Supp}(\mathscr F_y)=p_y^{-1}(\operatorname{Supp}(\mathscr F))
=\operatorname{Supp}(\mathscr F)\cap f^{-1}(y)$ en vertu de (I, 9.1.13.1) et
(I, 3.6.1), puisque $\mathscr F$ est de type fini~; en outre
$\operatorname{Supp}(\mathscr F)$ est fermé dans $X$ (0\textsubscript{I},
5.2.2), et si $Z$ est un sous-préschéma fermé de $X$ ayant pour espace
sous-jacent $\operatorname{Supp}(\mathscr F)$ (I, 5.2.1), et $j$ l'immersion
canonique $Z\to X$, $f\circ j$ est localement de type fini (1.3.4)~; cela
prouve que la première assertion entraîne la seconde. Pour prouver la
première assertion, remarquons d'abord que l'on peut supposer $X$ et $Y$
réduits (1.3.4, (vi)), autrement dit $Y$ intègre. Soit $\xi$ un point de
$f^{-1}(\eta)$~; on peut remplacer $X$ et $Y$ par des voisinages ouverts
affines de $\xi$ et $\eta$ respectivement, et par suite supposer que
$X=\operatorname{Spec}(B)$, $Y=\operatorname{Spec}(A)$, $B$ étant une
$A$-algèbre de type fini. En outre, si $Z$ est le sous-préschéma fermé réduit
de $X$ ayant $\overline{\{\xi\}}$ pour ensemble sous-jacent (I, 5.2.1), on
peut, comme ci-dessus, remplacer $X$ par $Z$, autrement dit supposer $B$
\emph{intègre} (I, 5.1.4). Comme le morphisme $f$ est alors dominant,
l'homomorphisme correspondant $\varphi:A\to B$ est injectif (I, 1.2.7)~;
donc le corollaire résulte finalement du lemme (1.8.5.1).
\end{proof}

